% Generated draft: requires mathematical and source review.
\documentclass[11pt]{article}
\usepackage[margin=25mm]{geometry}
\usepackage{fontspec}
\setmainfont{Latin Modern Roman}
\usepackage{amsmath,amssymb,mathtools,unicode-math}
\setmathfont{Latin Modern Math}
\usepackage{enumitem,needspace,xurl,hyperref}
\setlength{\parindent}{0pt}
\setlength{\parskip}{5pt}
\setlength{\emergencystretch}{3em}
\newcommand{\N}{\mathbb{N}}
\newcommand{\Z}{\mathbb{Z}}
\newcommand{\Q}{\mathbb{Q}}
\newcommand{\R}{\mathbb{R}}
\newcommand{\C}{\mathbb{C}}
\newcommand{\F}{\mathbb{F}}
\newcommand{\A}{\mathbb{A}}
\newcommand{\PP}{\mathbb{P}}
\newcommand{\E}{\mathbb{E}}
\newcommand{\eps}{\varepsilon}
\newcommand{\dd}{\,\mathrm{d}}
\begin{document}
\section*{UnsolvedMath ID: 30001970}
\textbf{Source problem code:} catalogue entry
\subsection{Definitions and mathematical statement}
The supplied record uses ``medieval mathematical algorithm'' for a historical algorithmic text and ``omitted steps'' for steps absent from its written presentation. Neither term is given a precise mathematical definition. ``Tacit formulation'' concerns knowledge left unstated; ``textual corruption'' is offered as one possible explanation for an omission, not as an established fact about any specified text.

The source asks: When a medieval mathematical algorithm omits steps, is the tacit formulation intentional and meant to transmit the knowledge clearly, or does it arise from some other cause such as textual corruption?

This is a historical, pedagogical, or methodological inquiry. The supplied record specifies no algorithm, manuscript, omitted step, intended readership, evidential criterion for intention, or mathematical assertion with defined domains and quantifiers. The possible explanations are not stated to be exhaustive or mutually exclusive. Consequently, a precise mathematical proposition cannot be recovered from this record without adding unsupported hypotheses. The source question is retained as a discussion prompt requiring contextual historical investigation.
\subsection{Short English statement}
The question asks whether missing steps in medieval algorithmic texts reflect an intentional way of communicating knowledge or another cause, such as textual corruption. No particular text or criterion for deciding between explanations is supplied, so the record does not determine a mathematical proposition.
\subsection{Sources}
\begin{itemize}
\item Local UnsolvedMath registry, supplied record with permanent ID \(30001970\), title \emph{Tacit Steps and Textual Corruption in Medieval Algorithms}; source problem code \texttt{catalogue entry}, in the Oberwolfach Reports Open Problems collection. The record gives the proposed year as 2012 and identifies no individual proposer.
\item \emph{Explicit Versus Tacit Knowledge in Mathematics} (2012), Oberwolfach Reports, pp.\ 131--210, DOI \texttt{10.4171/owr/2012/04}. Supplied source citation and literature-source link, labeled ``original Oberwolfach report'': \url{https://doi.org/10.4171/owr/2012/04}. No specific page or passage for this question is supplied.
\item Supplied literature assessment, checked on \texttt{2026-08-21}: the extracted question is a conceptual inquiry about authorial tacit formulation, not a resolved mathematical proposition. The supplied status is \texttt{open}. These are supplied dated curation, not independently verified findings; the background expressly calls for independent verification of time-sensitive status and cited claims.
\item Supplied statement-verification note: the workshop report treats the question as a historical, pedagogical, or methodological discussion prompt rather than an unresolved mathematical proposition. The statement status is \texttt{unrecoverable}; no separate clean statement is promoted, and the original record is retained for provenance. This description is supplied registry curation, not an independent examination of the report.
\end{itemize}
\end{document}
